\subsection{Open-source rigid-body simulators}
\label{sec:engine-comparison}

We compare the STEP/HHT simulation with native MuJoCo~\cite{mujoco}
and Drake~\cite{drake}. The two controlled problems are a sphere
approaching a plane and a collision of two translating spheres.
Each mobile body has mass 1 kg. The approaching sphere has initial
speed 0.45 m/s. Both problems start 18 mm outside a 12 mm contact
layer and use an undamped normal stiffness of 5000 N/m, with
gravity and friction set to zero. The analytical reference uses
the one-contact half sine with reduced mass 1 or 0.5 kg. The
two-sphere case also conserves center-of-mass velocity.

The layer is represented by offset collision surfaces. The plane
is raised by the layer thickness in the first problem. In the
second problem each collision radius grows by half that thickness.
The three simulators use these offsets in their own contact
representations. Our directional STEP probes evaluate the offset
surfaces, and the trajectory-wide clearance tests are reported
separately in Section~\ref{sec:continuous-experiments}.
MuJoCo \EngineMuJoCoVersion{} uses its native one-dimensional
contact and direct reference stiffness scaled by the contact
inverse mass. Drake \EngineDrakeVersion{} uses native continuous
point contact with zero Hunt--Crossley dissipation and geometry
stiffnesses of 10,000 N/m in series.
The maximum force-law residual on the computed states is
\EngineForceResidual{} N across the tested steps.

\begin{table}[!htbp]
\centering\small
\caption{Native rigid-body contact at a 1 ms nominal step. Position
errors use all recorded coordinates and the analytic trajectory.
Force errors use the same sample times and the analytic pulse.}
\label{tab:engine-error}
\begin{tabular}{lrrrr}
\toprule
 & \multicolumn{2}{c}{Position RMSE (\(\mu\)m)} &
 \multicolumn{2}{c}{Force RMSE (N)} \\
\cmidrule(lr){2-3}\cmidrule(lr){4-5}
Method & Sphere--plane & Two spheres & Sphere--plane & Two spheres \\
\midrule
\EngineBenchmarkRows
\bottomrule
\end{tabular}
\end{table}

\begin{figure}[!htbp]
\centering
\includegraphics[width=\linewidth]{\EngineFigurePath}
\caption{Rigid-body comparison with native MuJoCo and Drake.
All trajectories use the same isolated normal force law.
Position errors are evaluated against the analytical solution.
The force panel shows the sphere--plane pulse at 1 ms.}
\label{fig:engine-comparison}
\end{figure}

Table~\ref{tab:engine-error} and Figure~\ref{fig:engine-comparison}
show the effect of the integration method on the same force pulse.
At 1 ms, MuJoCo RK4 and Drake RK3 have smaller position errors
than HHT in both examples. The archived runs include all three
steps, seven integration timings per configuration, and every
scored state. The wider CAD experiments evaluate different costs,
including trim classification and interval exclusion.

\ifdefined\EngineTimingRows
\begin{table}[!htbp]
\centering\small
\caption{Median integration time for 160 steps at 1 ms in the
same Ubuntu WSL environment. Model setup and trajectory export
are outside the timed region.}
\label{tab:engine-time}
\begin{tabular}{lrr}
\toprule
Method & Sphere--plane (\(\mu\)s) & Two spheres (\(\mu\)s) \\
\midrule
\EngineTimingRows
\bottomrule
\end{tabular}
\end{table}
\else
The STEP/HHT timings use Windows, and both engine timings use Ubuntu
WSL. The integration timings are archived with their environments.
\fi
